\documentclass[11pt,a4paper]{article}

% Paquetes recomendados
\usepackage[utf8]{inputenc}
\usepackage[spanish]{babel}
\usepackage{amsmath, amsthm, amssymb}
\usepackage{geometry}
\usepackage{listings}
\usepackage{xcolor}
\usepackage{booktabs}
\usepackage{graphicx}

\geometry{margin=2.5cm}

% Definici�n de colores para c�digo
\definecolor{codegray}{rgb}{0.5,0.5,0.5}
\definecolor{backcolour}{rgb}{0.95,0.95,0.92}

\lstset{
    backgroundcolor=\color{backcolour},
    basicstyle=\ttfamily\footnotesize,
    breakatwhitespace=false,
    breaklines=true,
    captionpos=b,
    keepspaces=true,
    showspaces=false,
    showstringspaces=false,
    showtabs=false,
    tabsize=2
}

% Entornos de Teoremas
\newtheorem{theorem}{Teorema}

\title{Torre de Hanoi: Estrategia �ptima, algoritmo de resoluci�n y an�lisis matem�tico}
\author{M. en C. Diego Espejel}
\date{\today}

\begin{document}

\maketitle

\begin{abstract}
Este art�culo presenta un estudio sistem�tico sobre la resoluci�n �ptima del juego de la Torre de Hanoi en cualquiera de sus configuraciones. Se formalizan los conceptos fundamentales y se adopta una notaci�n homog�nea basada en principios de programaci�n. Se plantean dos teoremas centrales ---Irrelevancia de Elementos Superiores (IES) y Reducci�n de Elementos Finalizados (REF)--- que fundamentan la estrategia �ptima. Finalmente, se extiende el m�todo a configuraciones arbitrarias (no primitivas) mediante las transformaciones $Red1()$ y $Red2()$.
\end{abstract}

\section{Introducci�n}
El problema de las Torres de Hanoi, introducido por �douard Lucas en 1883, es un ejemplo paradigm�tico de recursividad \cite{cite:242}. Tradicionalmente se estudia la versi�n primitiva ($H_0$), pero este trabajo aborda la versi�n generalizada donde los discos pueden estar distribuidos arbitrariamente entre las columnas \cite{cite:237, 292}.

\section{Definiciones y Notaci�n}
Para un sistema de $n$ discos $H_n = \{1, 2, \dots, n\}$, se definen tres columnas: $col\_i$ (inicial), $col\_a$ (auxiliar) y $col\_f$ (final) \cite{cite:250, 298}.

\subsection{Notaci�n de Movimientos}
Cada acci�n se expresa como una llamada a funci�n $move(a, b)$, donde $a$ es el disco o semipila a mover y $b$ es el destino \cite{cite:301}.
\begin{itemize}
    \item $move(1, col\_f)$: mueve el disco 1 a la columna final \cite{cite:303}.
    \item $move(stack[r], col\_f)$: mueve una semipila de $r$ discos, equivalente a $2^r - 1$ movimientos \cite{cite:305, 426}.
\end{itemize}

\section{Fundamentos Te�ricos}

\begin{theorem}[Irrelevancia de Elementos Superiores - IES]
Toda secuencia de movimientos que opere exclusivamente sobre un subconjunto de discos inferiores $S_{inf} = \{1, \dots, m\}$ es independiente de la disposici�n de los discos en $S_{sup} = \{m+1, \dots, n\}$, siempre que estos �ltimos no se muevan \cite{cite:255, 313}.
\end{theorem}

\begin{theorem}[Reducci�n de Elementos Finalizados - REF]
Si los discos de $S_{sup} = \{m+1, \dots, n\}$ se encuentran ya posicionados en $col\_f$ respetando el orden, la resoluci�n �ptima es equivalente a resolver la instancia restringida a $S_{inf}$ \cite{cite:259, 316}.
\end{theorem}

\section{Resoluci�n de Torres Primitivas ($H_0$)}
La soluci�n �ptima para $n$ discos sigue una recurrencia fundamental:
\begin{equation}
    minMov(H_0[n]) = 2 \cdot minMov(H_0[n-1]) + 1 = 2^n - 1
\end{equation}
Esta secuencia se construye mediante tres fases: mover $n-1$ discos a $col\_a$, mover el disco $n$ a $col\_f$, y finalmente trasladar los $n-1$ discos de $col\_a$ a $col\_f$ \cite{cite:349, 351, 364}.

\section{Resoluci�n de Torres No Primitivas}
Para configuraciones arbitrarias, se aplican iterativamente los siguientes operadores \cite{cite:267, 428}:

\begin{enumerate}
    \item \textbf{Red2()}: Descarta los discos que ya ocupan su posici�n final en $col\_f$ \cite{cite:414}.
    \item \textbf{Red1()}: Agrupa semipilas v�lidas $\{1, \dots, r\}$ como una unidad $stack[r]$ \cite{cite:262, 422}.
    \item \textbf{Liberaci�n}: Se busca liberar el disco mayor restante $n$ para moverlo a $col\_f$ \cite{cite:321, 431}.
\end{enumerate}

\section{Algoritmo y Diagrama de Flujo}
El proceso de decisi�n completo se rige por el estado de las columnas y la disponibilidad de los discos para ser movidos. El algoritmo converge cuando todos los discos alcanzan la configuraci�n $H_f$ \cite{cite:460, 463}.

\begin{table}[h]
\centering
\begin{tabular}{@{}ll@{}}
\toprule
Variable/Funci�n & Descripci�n \\ \midrule
$libre(k)$ & El disco $k$ no tiene nada encima. \\
$stack[r]$ & Discos $\{1, \dots, r\}$ tratados como unidad. \\
$M(n)$ & Secuencia �ptima para $H_0[n]$. \\ \bottomrule
\end{tabular}
\caption{Funciones y s�mbolos definidos \cite{cite:298}.}
\end{table}

\end{document}